\clearpage
\VerdierPhysicalPageAnchor{0245}%
\begingroup\small\setlength{\parskip}{0pt}\setlength{\emergencystretch}{1.5em}
\setlength{\abovedisplayskip}{4pt plus 2pt minus 2pt}%
\setlength{\belowdisplayskip}{4pt plus 2pt minus 2pt}%

\VerdierSemanticAnchor{ch4d-4-6-3-b}
\noindent\textit{b) Soit
$Q\colon K(\mathcal A)\longrightarrow D(\mathcal A)$ le foncteur de
passage au quotient. Le complexe $I^\bullet$ est $Q$-libre à droite
(chap.~II, 2.3.3), i.e.\ pour tout complexe $Y$, le morphisme canonique :}
\[
\operatorname{Hom}_{K(\mathcal A)}(Y,I^\bullet)
\longrightarrow\operatorname{Hom}_{D(\mathcal A)}(Y,I^\bullet)
\]
\textit{est un isomorphisme.}

\VerdierSemanticAnchor{ch4d-4-6-3-c}
\noindent\textit{c) Soit :}
\[
H^*(I^\bullet)\colon\cdots\longrightarrow H^p(I^\bullet)
\longrightarrow H^{p+1}(I^\bullet)\longrightarrow\cdots
\]
\textit{le complexe dont les composants sont les $H^p(I^\bullet)$,
$p\in\mathbb Z$, et les différentielles sont nulles. Il existe un et un
seul isomorphisme dans $K(\mathcal A)$
$I^\bullet\longrightarrow H^*(I^\bullet)$ induisant l'identité sur les
objets de cohomologie.}

\VerdierSemanticAnchor{ch4d-4-6-3-d}
\noindent\textit{d) Les morphismes canoniques
$H^*(I^\bullet)\longrightarrow H^p(I^\bullet)[-p]$ définissent des
isomorphismes dans $K(\mathcal A)$ et dans $D(\mathcal A)$ :}
\[
H^*(I^\bullet)\xrightarrow{\sim}
\prod_{p\in\mathbb Z}H^p(I^\bullet)[-p].
\]

\VerdierSemanticAnchor{ch4d-4-6-3-proof-abcd}
Montrons que a) entraîne b). En effet, il suffit (chap.~II, 2.3.3) de
montrer que pour tout complexe acyclique $Y$,
$\operatorname{Hom}_{K(\mathcal A)}(Y,I^\bullet)=0$, ce qui résulte de a).
Montrons que a) entraîne c). En effet, il suffit de remplacer, dans
l'isomorphisme de a), le complexe $Y$ par $I^\bullet$ et le complexe
$I^\bullet$ par $H^*(I^\bullet)$. On voit de même que a) entraîne d) en
remplaçant, dans l'isomorphisme de a), le complexe $I^\bullet$ par le
complexe $H^*(I^\bullet)$. Démontrons a). Il résulte immédiatement des
définitions que le complexe $I^\bullet$ est isomorphe à un complexe du
type :

\VerdierSemanticAnchor{ch4d-4-6-3-proof-a-decomposition}
\[
\resizebox{0.98\linewidth}{!}{$
\displaystyle
\cdots\longrightarrow
J^{n-2}\oplus H^{n-1}\oplus J^{n-1}
\xrightarrow{\left(\begin{smallmatrix}
0&0&\mathrm{id}\\0&0&0\\0&0&0
\end{smallmatrix}\right)}
J^{n-1}\oplus H^n\oplus J^n
\xrightarrow{\left(\begin{smallmatrix}
0&0&\mathrm{id}\\0&0&0\\0&0&0
\end{smallmatrix}\right)}
J^n\oplus H^{n+1}\oplus J^{n+1}
\longrightarrow\cdots
$}.
\]

\VerdierSemanticAnchor{ch4d-4-6-3-proof-a-homotopy}
et, par suite, il existe un homotopisme entre $I^\bullet$ et le complexe
$H^*(I^\bullet)$. Pour démontrer a), on peut donc supposer que
$I^\bullet=H^*(I^\bullet)$. Il est clair qu'on a alors un isomorphisme :

\VerdierSemanticAnchor{ch4d-4-6-3-proof-a-hom}
\[
\operatorname{Hom}_{K(\mathcal A)}\bigl(Y,H^*(I^\bullet)\bigr)
\xrightarrow{\sim}
\prod_{p\in\mathbb Z}
\operatorname{Hom}_{K(\mathcal A)}
\bigl(Y,H^p(I^\bullet)[-p]\bigr).
\]
De plus, les objets $H^p(I^\bullet)$ étant injectifs, le morphisme
canonique :
\[
\operatorname{Hom}_{K(\mathcal A)}
\bigl(Y,H^p(I^\bullet)[-p]\bigr)
\longrightarrow
\operatorname{Hom}_{\mathcal A}\bigl(H^p(Y),H^p(I^\bullet)\bigr)
\]
est un isomorphisme.
\endgroup
